\begin{abstract}
Large language models are increasingly composed into agent loops that
plan, call tools, and resume the same task after each action.
These loops press a shared memory hierarchy harder than conventional
multi-turn chat, because they hold a growing key-value (KV) prefix
across tool waits and place many sessions on one static
random-access memory (SRAM) and high-bandwidth memory (HBM) pool,
so that eviction and hierarchical placement become a session-level
efficiency problem orthogonal to compute-mode optimization.
Existing proxies based on recency, timeout, or identity miss the
mechanism information of the loop and therefore treat a live wait
as a cold, discardable unit. To
address this problem, we present Unified Native Inter-turn Session
Orchestration Nexus (\unison), an event-driven near-memory scheduler beside the
memory hierarchy in which Survival-Penalty Eviction for Agent
Return-gap (\spear) and Tiering in Idle-window DMA Events (\tide)
share one live ranking. \spear selects who leaves from a gap average
and a turn-indexed hazard, while \tide spends the observed wait as a
direct memory access (DMA) budget for who sits in the fast tier. On
both coding and general-mission benchmarks with three model
families, totaling 1\,415 sessions and 33\,596 turns, the joint policy
is the best non-oracle entry on hit rate and average memory access
time (AMAT) on every trace, raising the hit rate by $0.3\%$ to
$23.1\%$ and reducing AMAT by $22\%$ to $51\%$, and lowering time to
first token (TTFT) by $58\%$ to $89\%$ on the long-horizon traces. A structural necessity analysis shows that the unified near-memory
design cannot be decomposed into independent IPs or realized in
software without re-introducing documented failure modes. The 28-nm
CMOS scheduling core occupies 0.169\,mm$^{2}$ at 13.6\,mW and
150\,MHz, a negligible overhead relative to the KV
hierarchy it manages, and reproduces the floating-point ranking at a
Kendall $\tau$ exceeding 0.998.
\end{abstract}

\begin{IEEEkeywords}
session KV residency, memory tiering, eviction, LLM agents,
near-memory scheduler
\end{IEEEkeywords}
